\begin{python0}
from solutions import *; clear()
\end{python0}
\sectionthree{Brute force search}

Computer Science, like Mathematics, seeks to find solutions to problems.
One way to look for solutions is through ``brute-force''search. This
means going over all possible candidates and seeing which ones are the
right solutions.

Let's try one. Suppose you want to find \textbf{integer} solution(s) to this equation:
\begin{center}
x$^{3}$ = 729
\end{center}
If x is a negative integer, x$^{3}$ is negative. So we are definitely only interested in positive integers. Furthermore x must be less than 729 (well ... a lot less!). So we can try x = 0, 1, 2, ..., 729 and see which one satisfies the above equation.

WHOA!!!

That sure looks suspiciously like a for-loop to me ...

\begin{ex}
Write a for-loop with an index variable running from
0 to 729 to find (and print ... of course!) solution(s) to the equation

\begin{center}
x$^{3}$ = 729
\end{center}
\end{ex}
Well \ldots the above example is kind of silly since we could have
computed the cube root of 729 and see if it's an integer
value \ldots but what if you have \textit{\textbf{two}} variables in the
equation? Suppose you want to solve for positive integers x and y
satisfying

\begin{center}
x$^{2}$ + y$^{2}$= 13$^{2}$
\end{center}

Well first of all each x and y must be at most 13, right? So all you
need to do is to check all possible 0, ..., 13 for x and 0, ..., 13 for
y. In other words you need to check all the following cases:

x = 0, y = 0\\
x = 0, y = 1\\
x = 0, y = 2\\
...\\
x = 0, y = 13\\
x = 1, y = 0\\
x = 1, y = 1\\
x = 1, y = 2\\
...\\
x = 1, y = 13\\
x = 2, y = 0\\
x = 2, y = 1\\
x = 2, y = 2\\
...\\
x = 2, y = 13\\
...\\
...\\
...\\
x = 13, y = 13\\

That sure looks like a double-nested for-loop to me!
Here's the program:
\begin{console}
for (int x = 0; x <= 13; ++x)
{   
    for (int y = 0; y <= 13; ++y)
    {
        if (x * x + y * y == 13 * 13)
        {  
           std::cout << x << ',' << y << '\n';
        }
    }
}
\end{console}
Now note that you see some ``repetitions''. I'm saying
that if x = a, y = b is a solution, then x = b, y = a is also a solution
because in the equation above, you can switch x and y. What if you do
not want to include such repetitions? One way would be to ensure x
<= y. In other words you try to get your program to run through
these values of x,y:

x = 0, y = 0\\
x = 0, y = 1\\
x = 0, y = 2\\
...\\
x = 0, y = 13\\
x = 1, y = 1 (note that y starts with 1 and not 0)\\
x = 1, y = 2\\
x = 1, y = 3\\
...\\
x = 1, y = 13\\
x = 2, y = 2 (note that y starts with 2 and not 0)\\
x = 2, y = 3\\
x = 2, y = 4\\
...\\
x = 2, y = 13\\
...\\
...\\
...\\
x = 13, y = 13\\
\begin{console}[commandchars=\~\@\$]
for (int x = 0; x <= 13; ++x)
{   
    for (int y = ~textbf@x$; y <= 13; ++y)
    {
        if (x * x + y * y == 13 * 13)
        {  
           std::cout << x << ',' << y << '\n';
        }
    }
}
\end{console}
Of course you realize that x$^2$ + y$^2$ = 13$^2$ gives the solution of the right-angle triangle with hypotenuse of length 13 and integer sides. Why not find all such triangles with hypotenuse from 1 to 100??? This means solving

\begin{center}
x$^2$ + y$^2$ = z$^2$
\end{center}

where 1 <= z <= 100.

\begin{console}[commandchars=\~\@\$]
~textbf@for (int z = 1; z <= 100; ++z)$
~textbf@{$   
    for (int x = 0; x <= z; ++x)
    {   
        for (int y = x; y <= z; ++y)
        {   
            if (x * x + y * y == z * z)
            {  
               std::cout << x << ','
                         << y << ','
                         << z << '\n';
            }
        }
    }
~textbf@}$
\end{console}

This prints the sides of all right angle triangles with integer sides
where the hypotenuse is between 1 to 100.

\begin{ex}
Can you improve the performance of the above program?
\end{ex}
\begin{ex}
Find all integer solutions to the equation
\begin{center}
x$^{4}$ + 2y$^{4}$ = 10001
\end{center}
\end{ex}
\begin{ex}
Here's something from wikipedia.org:
\begin{console}[commandchars=\~\@\$]
~textbf@1729$ is known as the ~textbf@Hardy-Ramanujan number$, after a
famous anecdote of the British mathematician
~href@http://en.wikipedia.org/wiki/G._H._Hardy$@G. H. Hardy$ regarding a
hospital visit to the Indian mathematician
~href@http://en.wikipedia.org/wiki/Srinivasa_Ramanujan$@Srinivasa
Ramanujan$. In Hardy's words:

~lq~lq I remember once going to see him when he was ill at
~href@http://en.wikipedia.org/wiki/Putney$@Putney$. I had ridden in taxi cab number 1729
and remarked that the number seemed to me rather a dull
one, and that I hoped it was not an unfavorable ~href@http://en.wikipedia.org/wiki/Omen$@omen$.
"No," he replied, "it is a very interesting number;
~textbf@it is the smallest number expressible as the sum of two$
~textbf@cubes in two different ways."$

The quotation is sometimes expressed using the
term "positive cubes", as the admission of negative
perfect cubes (the cube of a ~href@http://en.wikipedia.org/wiki/Negative_and_non-negative_numbers$@negative$ ~href@http://en.wikipedia.org/wiki/Integer$@integer$) gives
the smallest solution as ~href@http://en.wikipedia.org/wiki/91_%28number%29$@91$ (which is a factor of 1729):

~tab[3em]@91 = 6^3 + (-5)^3 =$
~tab[3em]@4^3 + 3^3$

Of course, equating "smallest" with "most negative", as
opposed to "closest to zero" gives rise to solutions
like -91, -189, -1729, and further negative numbers.
This ambiguity is eliminated by the term "positive cubes".

Numbers such as

~tab[3em]@1729 = 1^3 + 12^3 =$
~tab[3em]@9^3 + 10^3$

that are the smallest number that can be expressed as the
sum of two cubes in ~texttt@n$ distinct ways have been dubbed
~href@http://en.wikipedia.org/wiki/Taxicab_number$@taxicab numbers$. 1729 is the second taxicab number (the
first is 2 = 1^3 + 1^3). The number was also found in one
of Ramanujan's notebooks dated years before the incident.
\end{console}
Find all positive integer solutions to the equation
\begin{center}
x$^3$ + y$^3$ = 1729
\end{center}
(Hint: Write a double for-loop.)
\end{ex}
\begin{ex} The above exercise verified that there are exactly
two integer solutions. It does not verify that 1729 is the
``\textbf{smallest number expressible as the sum of two cubes in two
different ways}''. To do that you need to solve this:
\begin{center}
x$^3$ + y$^3$ = z
\end{center}
Of course you should print out x, y, and z. Look at all the solutions
and find the value of z with exactly two integer solutions. Is it 1729?
(Or is Ramanujan totally wrong all these years ...)
\end{ex}
\begin{ex}
Write a program that prints all positive fractions m/n from 0
to 5 where m and n are positive and at most 10. The fractions need not
be reduced or in any particular order or unique.
\end{ex}

